\section{定位：A 股长多研究与市场中性边界}

风格因子的公开学术与机构语境，通常建立在大样本、规则化再平衡、点时数据、成本控制和（在许多市场）多空组合之上。中国个人的 A 股研究可借用质量、价值、动量和投资等特征来组织公开资料；稳定空头腿、融资、借券和对冲能力缺失时，市场中性、套利和纯因子收益等表述均缺少结构基础。本章只定义一个\textbf{长多候选排序与否决协议}；它不输出代码、权重或非零持仓。

\section{先定义可研究宇宙，再计算分数}

候选宇宙必须在每个再平衡日冻结；退市名单、财务字段和行业分类均须采用当时版本。任何研究日应保存证券身份、交易状态、上市日期、板块/行业分类、财务报告可得时间、价格/成交字段、风险警示标志和数据版本。下列过滤器构成研究协议中的最低字段要求，不产生对某个股票、行业或产品的推荐：

\begin{exhibitbox}[Exhibit 7：A 股因子候选的准入与否决门]
\small
\begin{tabularx}{\exhibitboxwidth}{L{2.6cm}X L{3.65cm}}
\toprule
\textbf{门槛类别} & \textbf{纸面研究必须回答的问题} & \textbf{否决 / 降级逻辑} \\
\midrule
身份与可得时间 & 是否为当期可验证的境内上市证券？财务字段在截点前是否真的已公开？ & 不能以今日修订的年报或分类回填过去。 \\
交易状态 & 是否处于停牌、风险警示、连续价格限制、特殊上市阶段或其他异常状态？ & 不能把“触发信号”当成一定可以处理；状态不明即剔除。 \\
流动性与容量 & 价格、成交额、换手、盘口/交易成本资料是否足以支持研究假设？ & 单日成交额、热度或榜单不能单独证明容量；缺字段不排序。 \\
财务可比性 & 质量、价值和投资字段的会计口径、报告期、币种、非经常项处理是否一致？ & 合并/拆分、异常利润、审计保留、字段修订均须显式处理。 \\
集中度 & 行业、主题、供应链、客户和共同宏观因子的重叠是否被识别？ & 多只代码不等于分散；不可解释重叠则不能进入纸面篮子。 \\
\bottomrule
\end{tabularx}
\sourcenote{S10、S13；本报告的点时数据和市场制度研究协议。}
\end{exhibitbox}

\section{冻结分数：用于复核，不用于追逐}

在通过准入门之后，研究协议可以计算一个横截面复合分数：

\[
F_{j,t}=0.30Q_{j,t}+0.25V_{j,t}+0.25M_{j,t}+0.20I_{j,t},
\]

其中 \(Q\) 代表质量，\(V\) 代表价值，\(M\) 代表动量，\(I\) 代表投资/资本配置。每一子分数均应先在可比行业或其他预先固定分组中标准化，再记录原始字段、转换方法、截点和数据版本。复合分数用于把“好公司”“便宜”“趋势强”等模糊词拆成可被复核的字段；公司价值、收益概率和买入理由需要独立证据。

\begin{dashboardbox}[四类因子的纸面定义与常见误用]
\begin{tabularx}{\textwidth}{L{2.0cm}X L{4.2cm}}
\toprule
\textbf{因子} & \textbf{可审计的研究问题} & \textbf{不能被偷换成} \\
\midrule
质量 \(Q\) & 盈利质量、现金转换、资本回报、资产负债表和治理是否在同一报告时点下可解释？ & 高估值成长叙事、一次性利润或“知名公司”标签。 \\
价值 \(V\) & 市场价格相对可解释的利润、现金流、资产或经营价值桥处于何种区间？ & 低市盈率/市净率必然便宜；周期高点利润或低质量资产的折价。 \\
动量 \(M\) & 价格/盈利修正/经营数据的延续在何时可得，是否受停牌、涨跌停和流动性扭曲？ & 追逐涨停、媒体热度或将短期价格动量写成基本面验证。 \\
投资 \(I\) & 资本开支、并购、营运资本和扩张是否能以公开资料解释其未来回报与现金需求？ & 低资本开支永远更优，或用“扩产”叙事自动给予估值信用。 \\
\bottomrule
\end{tabularx}
\end{dashboardbox}

\section{固定再平衡节奏和约束，避免把研究变成随意调仓}

本备忘录记录一个供未来审计使用的测试围栏：在月度或季度的预先固定评估点冻结数据、重新计算分数、检查行业和流动性约束、再记录“可研究 / 否决”的变化。为使围栏可被检验，协议保留以下\emph{非执行参数}：最多 20 个纸面候选、单一名称上限 \(6\%\)、单一行业上限 \(25\%\)、因子研究袖套总上限 \(40\%\)。这些数字仅供审计；当前纸面候选权重仍全部为 \(0\%\)。

若未来对该围栏进行回放，必须同时报告：全样本和留出样本、退市/暂停交易处理、财报修订、手续费与滑点假设、无法成交日、行业集中度、最大回撤、因子暴露变化和与简单基准的可比定义。只报告年化收益、胜率或夏普比率不足以说明因子框架有效。

\section{A 股结构如何改变“止损”与“再平衡”}

A 股的 T+1、停复牌和价格限制意味着价格阈值不能被假设为即时成交的订单。因子研究的首要风险控制应发生在进入前：限制单名和行业集中、排除交易状态不清的对象、把流动性不足视为否决、并在出现治理/审计/监管或核心基本面变化时触发重新承保。价格变动本身只能触发研究复核，不能替代证据。

\begin{riskbox}[因子框架的反直觉结论]
如果一个“因子策略”需要频繁临时更换指标、忽略财务发布日期、跳过停牌与价格限制、或把长多组合称为市场中性，研究应拒绝该结果。因子的用途在于约束主观性；它不能为主观追逐提供新名称。
\end{riskbox}
